\section{Detection Head}

\subsection{Idea generale}

La \textit{detection head} è il componente della rete che riceve le feature
estratte dal backbone e dalla FPN e le trasforma nelle predizioni necessarie
per l'object detection. In generale, una head può essere suddivisa in più
branch, ognuna specializzata in un compito differente, ad esempio la
classificazione dell'oggetto e la regressione della sua bounding box.

L'idea generale è quindi:

$$
\text{Feature Map}
\rightarrow
\text{Detection Head}
\rightarrow
\begin{cases}
\text{Classification} \\
\text{Localization}
\end{cases}
$$

Nel caso dei detector anchor-free, la predizione viene effettuata direttamente
su ogni posizione spaziale della feature map invece di utilizzare bounding box
predefinite come riferimento.

\subsection{RetinaNet}

Un esempio importante è RetinaNet. La rete utilizza due subnet separate,
una per la classificazione e una per la regressione delle bounding box.
Entrambe ricevono una feature map della FPN e sono costituite da più
convoluzioni $3\times3$, seguite da un layer finale di predizione.

La classification subnet produce, per ogni posizione spaziale, le probabilità
relative alle classi e agli anchor presenti in quella posizione, mentre la box
regression subnet produce i parametri necessari per correggere le coordinate
degli anchor.\footnote{Tsung-Yi Lin et al., \textit{Focal Loss for Dense Object
Detection}, ICCV 2017.}

La struttura può essere schematizzata come:

$$
\text{FPN feature}
\rightarrow
\begin{cases}
\text{Classification subnet} \\
\text{Regression subnet}
\end{cases}
$$

In RetinaNet le due subnet hanno la stessa struttura generale, ma utilizzano
parametri separati. La classification subnet termina con $KA$ output per
posizione, mentre la regression subnet termina con $4A$ output, dove $K$ è il
numero di classi e $A$ il numero di anchor per posizione.

\subsection{FCOS}

FCOS mantiene la separazione tra classificazione e regressione, ma elimina il
riferimento agli anchor. Ogni posizione della feature map viene trattata come
una possibile location dell'oggetto e la regressione predice direttamente le
distanze dalla posizione ai quattro lati della bounding box:

$$
(l,t,r,b).
$$

La detection head di FCOS contiene quindi una branch per la classificazione e
una branch per la regressione. Il paper descrive per entrambe una sequenza di
convolutional layers applicata alle feature della FPN. In aggiunta viene
introdotta una terza branch, la \textit{center-ness branch}, che predice quanto
una location sia vicina al centro dell'oggetto e serve a ridurre l'influenza
delle predizioni di bassa qualità.\footnote{Zhi Tian et al.,
\textit{FCOS: Fully Convolutional One-Stage Object Detection}, ICCV 2019.}

Lo schema concettuale è quindi:

$$
\text{FPN feature}
\rightarrow
\begin{cases}
\text{Classification} \\
\text{Regression} \\
\text{Center-ness}
\end{cases}
$$

La classificazione determina \textit{che cosa} è presente in una determinata
location, la regressione determina \textit{dove} si trova la bounding box,
mentre la center-ness fornisce un'indicazione della qualità della location
utilizzata per effettuare la predizione.

\subsection{Detection head del paper sulla pneumonia}

Il paper utilizzato come riferimento per questo progetto descrive una head più
semplice. Ogni livello della feature pyramid viene elaborato inizialmente da
una convoluzione $3\times3$ e successivamente da due convoluzioni $1\times1$
che producono due predizioni distinte, indicate come \textit{CENTER} e
\textit{SCALE}.

$$
\text{FPN feature}
\rightarrow
3\times3
\rightarrow
\begin{cases}
1\times1 \rightarrow \text{CENTER}\\
1\times1 \rightarrow \text{SCALE}
\end{cases}
$$

Il branch CENTER viene utilizzato per individuare le location associate al
centro degli oggetti, mentre il branch SCALE viene utilizzato per predire la
scala della bounding box.\footnote{Linghua Wu et al.,
\textit{Pneumonia detection based on RSNA dataset and anchor-free deep learning
detector}, Scientific Reports, 2024.}

\subsection{Detection head del progetto}

La detection head utilizzata nel progetto è più vicina a quella di FCOS che a
quella descritta nel paper sulla pneumonia.

Ogni livello della FPN $P_3,\ldots,P_7$ viene passato a una
\textit{DetectionHead}. La head è divisa in due tower principali:

$$
\text{Feature}
\rightarrow
\begin{cases}
\text{Classification tower}\\
\text{Regression tower}
\end{cases}
$$

La classification tower è costituita, nel progetto, da:

$$
3\times3\ \text{Conv}
\rightarrow
\text{GroupNorm}
\rightarrow
\text{ReLU}
\rightarrow
1\times1\ \text{Conv}.
$$

La regression tower utilizza una struttura analoga:

$$
3\times3\ \text{Conv}
\rightarrow
\text{GroupNorm}
\rightarrow
\text{ReLU}.
$$

Dalla regression tower vengono poi prodotte due predizioni separate:

$$
\text{Regression tower}
\rightarrow
\begin{cases}
1\times1 \rightarrow (l,t,r,b)\\
1\times1 \rightarrow \text{centerness}
\end{cases}
$$

Pertanto, per ogni location di ogni livello della FPN, il modello produce:

$$
\text{classification},\qquad
(l,t,r,b),\qquad
\text{centerness}.
$$

Nel progetto la classificazione produce attualmente un solo valore per
location, perché viene utilizzata una sola classe target, la pneumonia.

\subsection{Differenze rispetto ai paper}

La struttura del progetto non coincide esattamente con nessuno dei due paper.

Rispetto al paper sulla pneumonia, il progetto aggiunge una
\textit{classification branch} esplicita e utilizza una regressione
$LTRB$, più vicina alla formulazione FCOS, invece della distinzione CENTER/SCALE
descritta nel paper.

Rispetto a FCOS, la struttura è semplificata: il paper originale utilizza
quattro convolutional layers nelle branch di classificazione e regressione,
mentre il progetto utilizza una sola convoluzione $3\times3$ per ciascuna
tower, seguita da Group Normalization e ReLU.

Inoltre, FCOS condivide le head tra i diversi livelli della feature pyramid,
mentre nel progetto vengono create cinque istanze separate di
\texttt{DetectionHead}, una per ciascun livello $P_3,\ldots,P_7$. Il paper
FCOS descrive infatti la condivisione delle head tra i diversi livelli, pur
introducendo modifiche specifiche per gestire i differenti range di scala.

Nel complesso, quindi, la detection head del progetto può essere considerata
una \textit{FCOS-like detection head}: ne mantiene la formulazione
anchor-free, la classificazione per location, la regressione $LTRB$ e la
center-ness branch, ma utilizza una struttura più semplice e alcune scelte
architetturali specifiche del progetto.
\subsubsection{Group Normalization}

La \textit{Group Normalization} (GN) è una tecnica di normalizzazione delle
feature map che divide i canali in gruppi e normalizza i valori all'interno di
ciascun gruppo.

A differenza della Batch Normalization, la Group Normalization non utilizza le
statistiche dell'intero batch. Questo la rende meno dipendente dalla
dimensione del batch e quindi particolarmente utile quando si lavora con batch
molto piccoli, situazione comune nell'object detection a causa dell'elevato
costo computazionale delle immagini.

Nel progetto la Group Normalization viene utilizzata all'interno delle
classification e regression tower della detection head:

\[
\text{Conv}
\rightarrow
\text{GroupNorm}
\rightarrow
\text{ReLU}.
\]

Il suo scopo è stabilizzare le attivazioni e rendere l'addestramento più
robusto, senza richiedere batch grandi.